\section{The alternate collapse and Mayer--Vietoris route}

The nearby-cycle/Leray proof in \cref{sec:integral-leray} gives a verified
independent route to the topology of $X$.  This section keeps the manuscript's independent
cellular-collapse/Mayer--Vietoris route under audit.  The integer algebra in
\cref{prop:terminal-group,prop:terminal-degree-two} is complete, but the
geometric realization of the printed homology matrices is not yet needed or
fully validated.  Write the relevant decomposition as
$X=U\cup N$, with $N=\bigcup_jN_j$ and $U\cap N$ the disjoint union of
oriented collars.  The Mayer--Vietoris map is
\[
\delta_q=(i_{U*},-i_{N*}):
H_q(U\cap N;\mathbb Z)\longrightarrow
H_q(U;\mathbb Z)\oplus H_q(N;\mathbb Z).
\tag{4.1}
\]
The assertion requiring geometric proof is not the Smith form of a printed
matrix but
\[
\alpha_q=[\delta_q]_{\mathcal B_q}
\tag{4.2}
\]
for bases $\mathcal B_q$ constructed from actual primitive oriented cycles.

\subsection{A complete abstract cusp-retraction lemma}

\begin{proposition}[equivariant topological spine]\label{prop:topological-spine}
\Verified\ Let $A$ be a contractible manifold with boundary and
$A_0\subset\partial A$ a contractible union of boundary components.  Let a
group $\Gamma$ act freely and properly discontinuously on $(A,A_0)$.  Assume
$Q=A/\Gamma$ and $H=A_0/\Gamma$ are manifolds, $Q$ is compact,
$H\hookrightarrow Q$ is collared, and the two quotient maps are connected
coverings with deck group $\Gamma$.  If the inclusion identifies the two
deck groups by the identity automorphism of $\Gamma$, then
$H\hookrightarrow Q$ is a strong deformation retract.  The retraction lifts
uniquely from the identity lift to a $\Gamma$-equivariant strong deformation
retraction of $A$ onto $A_0$.
\end{proposition}

\begin{proof}
Because $A$ and $A_0$ are contractible universal covers,
$\pi_n(Q)=\pi_n(H)=0$ for $n\ge2$, while
$\pi_1(Q)\cong\Gamma\cong\pi_1(H)$.  By hypothesis the inclusion induces the
identity on $\Gamma$, so it is a weak homotopy equivalence.  Manifolds have
the homotopy type of CW complexes; Whitehead's theorem
\cite[Theorem~4.5]{Hatcher2002} therefore makes the inclusion a homotopy
equivalence.

The collar makes $(Q,H)$ an NDR pair, hence the inclusion is a cofibration.
We record the standard cofibration argument.  Choose a homotopy inverse
$r_0:Q\to H$.  Since $r_0|_H\simeq\operatorname{id}_H$, the homotopy
extension property modifies $r_0$ through maps $Q\to H$ to a retraction
$r$ satisfying $r|_H=\operatorname{id}_H$.  The maps
$\operatorname{id}_Q$ and $i\circ r$ are homotopic and agree on $H$.
Applying the homotopy extension property once more makes that homotopy
stationary on $H$.  It is therefore a strong deformation retraction.

Let $P:A\to Q$ be the covering.  Lift the deformation homotopy from the
initial lift $\operatorname{id}_A$.  Since the base homotopy fixes $H$ and
the initial lift fixes $A_0$, uniqueness of path lifting makes the lifted
homotopy fix $A_0$.  For $\gamma\in\Gamma$, both
$\widetilde H_t\circ\gamma$ and $\gamma\circ\widetilde H_t$ lift the same
map and agree at $t=0$; uniqueness gives equality for every $t$.  The lift is
therefore $\Gamma$-equivariant and ends in $A_0$.
\end{proof}

This proves the formal retraction step once its hypotheses are established.
The manifold, freeness, proper discontinuity, and quotient assertions for the
infinite toric space are now proved in
\cref{thm:candidate-cusp-geometry}; the candidate-specific collar and
contractibility assertions used by this alternate route remain to be checked.

\subsection{The remaining candidate-specific obligations}

\begin{obligation}[C1: finite collapse model]
\label{obl:finite-model}
\Open\ The analytic quotient prerequisites are discharged by
\cref{thm:candidate-cusp-geometry}.  For the alternate collapse proof,
construct the claimed finite quotient CW structure, verify its attaching maps
on all common faces, and check the collar, contractibility, and identity deck
group hypotheses needed by \cref{prop:topological-spine}.
\end{obligation}

\begin{obligation}[C2--C3: collapse and specialization]
\label{obl:cusp-collapse}
\Open\ Construct the collapse on each toric stratum, prove compatibility on
all face inclusions and $\bar\Lambda$-equivariance, and compute its map on
integral cellular chains with oriented primitive bases.  Then prove that this
topological collapse map agrees with the nearby-cycle specialization
\[
\operatorname{sp}^q:H^q(W;\mathbb Z)\longrightarrow H^q(F;\mathbb Z)
\].
The image, torsion, and saturation of specialization itself are proved in
\cref{thm:p0-integral-specialization} for the reconstructed cusp model.
What remains for this alternate route is the
model-specific equality of that map with the cellular collapse and the
oriented chain matrix.  SGA~7 supplies formal nearby cycles
\cite{SGA7II1973}; it does not supply that equality.
\end{obligation}

\begin{obligation}[V1: based van Kampen inclusions]
\label{obl:vankampen}
\Verified\ The actual overlap coordinate maps, with the zero section as
base locus and clockwise meridians, show that the
finite fillings impose $\widehat\rho_j^{m_j}=t_{v_j}$ with the same sign,
the cusp filling kills $\Lambda_{\rm tor}$ and its toric meridian, and the
tautological gluing has zero translation drift.  Thus the abstract terminal
presentation is the actual $\pi_1(X)$.  The deck generators, common clockwise
convention, zero-section disc, normal closure, boundary inclusion
$\Lambda\to\Lambda/\Lambda_{\rm tor}$,
$\widetilde a_0\mapsto1$, and absence of extra whisker drift are proved in
\cref{prop:geometric-pi1}.
\end{obligation}

\begin{obligation}[MV1 and the remaining comparison]
\label{obl:leray}
\Open\ Prove (4.2) for every degree used in the homology calculation.
The integral Leray transgression and its candidate sign are already derived
in \cref{lem:first-transgression}; the higher transgressions are
\cref{lem:higher-transgressions}.  What remains is a commuting comparison
diagram between those maps and (4.1), listing every
nearby-cycle, meridian, transfer, and basis identification shared by the two
routes.  Agreement of their terminal integers is not itself such a diagram.
\end{obligation}

\begin{test}[immediate topological falsifiers]
The alternate collapse certificate fails if any actual map has a
nonprimitive image where
the printed matrix assumes a primitive one, if specialization contributes an
omitted torsion class or nonunit index, if a meridian sign changes $p=-1$ to
a nonunit, or if the geometric Leray transgression differs from
$\pm p\omega$ by a nonunit factor.
\end{test}

\subsection{Complete conditional recognition theorem}

\begin{lemma}[cohomological closure in dimension six]
\label{lem:cohomological-closure}
\Verified\ Let $X$ be a smooth, closed, connected, oriented six-manifold.
If
\[
\pi_1(X)=0,
\qquad H^2(X;\mathbb Z)=H^3(X;\mathbb Z)=0,
\tag{4.3}
\]
then $H_k(X;\mathbb Z)$ is $\mathbb Z$ for $k=0,6$ and zero for
$1\le k\le5$.
\end{lemma}

\begin{proof}
The Hurewicz map in degree one identifies
$H_1(X;\mathbb Z)=\pi_1(X)^{\mathrm{ab}}=0$.  The universal coefficient
sequence
\[
0\longrightarrow\operatorname{Ext}(H_1(X),\mathbb Z)
\longrightarrow H^2(X;\mathbb Z)
\longrightarrow\operatorname{Hom}(H_2(X),\mathbb Z)
\longrightarrow0
\]
shows that $\operatorname{Hom}(H_2(X),\mathbb Z)=0$, so the finitely
generated group $H_2(X)$ is torsion.  The next universal coefficient
sequence contains an injection
\[
\operatorname{Ext}(H_2(X),\mathbb Z)
\hookrightarrow H^3(X;\mathbb Z)=0.
\]
For a finitely generated torsion group
$A=\bigoplus_i\mathbb Z/n_i$, one has
$\operatorname{Ext}(A,\mathbb Z)=\bigoplus_i\mathbb Z/n_i$; hence
$H_2(X)=0$.  Poincar\'e duality gives
\[
H_3(X)\simeq H^3(X)=0,\quad
H_4(X)\simeq H^2(X)=0,\quad
H_5(X)\simeq H^1(X)=0.
\]
Finally connectedness and orientation give $H_0(X)=H_6(X)=\mathbb Z$.
\end{proof}

\begin{theorem}[sphere recognition from the exact inputs]
\label{thm:sphere-recognition}
\Conditional\ Let $X$ be a smooth, closed, connected six-manifold.  If
\[
\pi_1(X)=0,\qquad
H_k(X;\mathbb Z)\cong
\begin{cases}
\mathbb Z,&k=0,6,\\
0,&1\le k\le5,
\end{cases}
\tag{4.4}
\]
then $X$ is diffeomorphic to the standard $S^6$.
\end{theorem}

\begin{proof}
A closed smooth manifold has the homotopy type of a finite CW complex.
We first prove successively that $\pi_i(X)=0$ for $1\le i\le5$.
The case $i=1$ is assumed.  Suppose $2\le n\le5$ and
$\pi_i(X)=0$ for $i<n$.  Then $X$ is $(n-1)$-connected, so the Hurewicz
homomorphism
\[
h_n:\pi_n(X)\longrightarrow H_n(X;\mathbb Z)
\]
is an isomorphism by the Hurewicz theorem
\cite[Theorem~4.32]{Hatcher2002}.  The target vanishes by (4.4), hence
$\pi_n(X)=0$.  Induction proves the assertion.

Now $X$ is $5$-connected.  Hurewicz in degree six gives
$\pi_6(X)\cong H_6(X;\mathbb Z)\cong\mathbb Z$.  Choose a map
$g:S^6\to X$ whose homotopy class maps to a chosen orientation generator.
It has degree one.  Hence $g_*$ is an isomorphism in degrees zero and six;
in every other degree both integral homology groups vanish.  Thus $g$ is an
integral homology equivalence.  Both spaces are simply connected CW
complexes, so the homological Whitehead theorem
\cite[Corollary~4.33]{Hatcher2002} makes $g$ a homotopy equivalence.
Therefore $X$ is a smooth homotopy $6$-sphere.

Smale's generalized Poincar\'e theorem in dimensions greater than four
\cite[Theorem~A]{Smale1961} identifies the underlying topological manifold
with $S^6$.  Smooth homotopy $6$-spheres up to oriented
$h$-cobordism form the group $\Theta_6$; the opening table of
Kervaire--Milnor gives $|\Theta_6|=1$
\cite{KervaireMilnor1963}.  Hence there is no exotic smooth structure in
dimension six, and $X$ is diffeomorphic to the standard sphere.
\end{proof}

\begin{corollary}[abstract application boundary]\label{cor:conditional-s6}
\Conditional\ Suppose V1 identifies the actual based van Kampen quotient
with $G(\ell_0,\ell_1,\ell_2)$ of
\cref{prop:terminal-group}, and suppose an independently constructed
MV1/L1 comparison proves
\[
H^2(X;\mathbb Z)\cong H^3(X;\mathbb Z)
\cong\mathbb Z/|p|\mathbb Z,
\qquad p=12\ell_0-4\ell_1-3\ell_2.
\tag{4.5}
\]
If $|p|=1$, then $X\cong_{\mathrm{diff}}S^6$.
\end{corollary}

\begin{proof}
For $|p|=1$, 
\cref{prop:terminal-group} and the assumed geometric identification give
$\pi_1(X)=0$, while (4.5) gives $H^2(X)=H^3(X)=0$.
\Cref{lem:cohomological-closure} proves (4.4), and
\cref{thm:sphere-recognition} gives the claimed diffeomorphism.
\end{proof}

The corollary is an abstract implication.  For the actual candidate, the
intrinsic nearby-cycle/Leray and based van Kampen route now proves the
hypotheses used by \cref{cor:candidate-sphere}.  C1--C3, MV1, and L1 remain
only the independent collapse/Mayer--Vietoris certificate; the additional open
work in this section concerns that alternate certificate.
